\documentclass[10pt,a4paper]{amsart}
\usepackage[margin=19mm]{geometry}
\usepackage{amsmath,amssymb,mathtools}
\usepackage[scr=rsfs,cal=euler]{mathalfa}
\usepackage{xfrac}
\usepackage[unicode,hidelinks,bookmarks=false]{hyperref}
\newcommand{\ie}{i.e.\ }
\newcommand{\cf}{cf.\ }
\newcommand{\resp}{respectively\ }
\newcommand{\Subsection}{\subsection}
\providecommand{\nobreakdash}{\nobreak\hbox{-}}
\setlength{\emergencystretch}{2em}
\multlinegap0pt
\usepackage[shortlabels]{enumitem}
\usepackage{amssymb}
\usepackage{mathtools}
\newtheorem{lemma}{Lemma}
\newtheorem{theorem}{Theorem}
\newtheorem{proposition}{Proposition}
\usepackage{cleveref}
\crefname{lemma}{Lemma}{Lemmas}
\crefname{theorem}{Theorem}{Theorems}
\crefname{proposition}{Proposition}{Propositions}
\newcommand{\BVloc}{\operatorname{BV_{loc}}}
\newcommand{\Law}{\operatorname{Law}}
\newcommand{\MC}{\operatorname{MC}}
\newcommand{\Pcal}{\mathcal{P}}
\newcommand{\R}{\mathbb{R}}
\newcommand{\W}{W}
\newcommand{\ip}[2]{\left\langle #1,#2\right\rangle}
\newcommand{\mtwo}{m_2}
\newcommand{\qmu}{\mathsf q_{\mu}}
\newcommand{\tr}{\operatorname{tr}}
\title{Couplings Farthest from the Independent Gaussian}
\author{Stefan Schrott}
\date{Source: arXiv:2609.10467v1 (CC BY 4.0)}
\begin{document}
\maketitle

\noindent\emph{Source and licence.} Couplings Farthest from the Independent Gaussian. Version arXiv:2609.10467v1 (CC BY 4.0), distributed by arXiv under the Creative Commons Attribution 4.0 International licence, http://creativecommons.org/licenses/by/4.0/, as recorded in the arXiv licence metadata for this exact version and shown on the abstract page https://arxiv.org/abs/2609.10467v1. Full licence notice: \texttt{licenses/arxiv-2609.10467-CC-BY-4.0.txt}.\par\smallskip

\noindent\textit{Editorial scope.} Counted unit: the article's theorem thm:general (source0.tex lines 607--643). The article's complete original proof of it is delivered with it (source0.tex lines 1019--1212, one \texttt{proof} environment of 194 lines, which the article places away from the statement). No mathematics is written, repaired or re-derived: the statement and the proof are the article's own lines, in the article's own notation, and no retained line is substituted or re-flowed.  Retained uncounted as the source-local context the proof needs: lines 455--465; lines 687--699; lines 848--861; lines 912--946.  Nothing was omitted from any retained span, so the package carries no omission record.  No retained line contains a citation, so the wrapper carries no bibliography and no bibliography entry.  No retained line contains a figure environment, an image-inclusion command or drawing code, so no image is delivered and the package declares no asset dependency.  Editorial formatting only, in addition: an AMS \texttt{amsart} preamble that restates the article's own declarations and macros used by the retained lines, namely the environments \texttt{lemma}, \texttt{theorem}, \texttt{proposition} on separate counters, together with the standard packages those lines need.  The retained environments print in this document as Lemma 1, Theorem 1, Proposition 1 in order of appearance, whereas the article prints the same items with the numbering of its own full document.  The complete native source of the article as distributed in the arXiv e-print for this exact version is bundled unchanged as \texttt{sources/209919/source0.tex}.
\begin{lemma}[Gaussian integration by parts for convex functions]\label{lem:GaussPartInt}
Let $\varphi \in L^1(\gamma_n)$ be convex and assume that
$\nabla\varphi\in L^2(\gamma_n;\mathbb R^n)$. Then
\[
    \int_{\mathbb R^n} \rho_n(x)\,d(\Delta\varphi)(x)
    =
    \int_{\mathbb R^n}
        \langle x,\nabla\varphi(x)\rangle\,d\gamma_n(x)
\]
and in particular, both sides of the equality are finite. 
\end{lemma}

\begin{proposition}\label{prop:ineq_component}
Let $\mu \in \Pcal_2(\R)$ and $u\in \BVloc(\R^n)$ satisfy
$u_\#\gamma_n=\mu$. Then
\begin{equation}\label{eq:scalar-bound}
  \int_{\R^n} \rho_n \,  d|Du| \ge \MC(\gamma,\mu).
\end{equation}
Equality holds if and only if there is a
unit vector $a\in\R^n$ such that
\begin{equation}\label{eq:scalar-equality}
  u(x)=\qmu(\ip{a}{x})
  \qquad\text{for $\gamma_n$-almost every $x$.}
\end{equation}
\end{proposition}

\begin{lemma}\label{lem:matrix}
Let $A\in\mathbb R^{n\times n}$ be symmetric and positive semidefinite
and write $e_i$ for the $i$-th unit vector. Then
\[
  \frac{1}{\sqrt n}\sum_{i=1}^n |Ae_i|
  \leq
  \operatorname{tr}(A).
\]
Equality holds if and only if
\[
  A=c\,\varepsilon\varepsilon^\top
\]
for some $c\geq0$ and some $\varepsilon\in\{-1,1\}^n$.
\end{lemma}

\begin{lemma}\label{lem:signed-gradient}
Let $\mu\in\Pcal_2(\R)$ and let
$\varepsilon\in\{-1,1\}^n$. Define
$T_\varepsilon\colon\R^n\to\R^n$ by
\[
  T_{\varepsilon,i}(x)
  :=
  \qmu\left(
    \frac{\varepsilon_i}{\sqrt n}
    \langle\varepsilon,x\rangle
  \right),
  \qquad i\in\{1,\ldots,n\},
\]
and define 
\[
  \varphi_\varepsilon(x)
  :=
  m_\mu\sum_{i=1}^n x_i
  +
  \sqrt n\,h_\mu\left(
    \frac{\langle\varepsilon,x\rangle}{\sqrt n}
  \right),
  \qquad
  h_\mu(s)
  :=
  \int_0^s\bigl(\qmu(r)-m_\mu\bigr)\,dr.
\]
Then $\varphi_\varepsilon$ is convex and the following statements hold:
\begin{itemize}
    \item If $\varepsilon=(1,\ldots,1)$ or $\varepsilon=(-1,\ldots,-1)$, then $T_\varepsilon=\nabla\varphi_{(1,\ldots,1)}$.
    \item If $\varepsilon$ contains both signs, then $T_\varepsilon$ is the
gradient of a convex function if and only if $\mu$ is symmetric about
$m_\mu$. In this case, $T_\varepsilon=\nabla\varphi_\varepsilon$.
\end{itemize}
\end{lemma}

\begin{theorem}\label{thm:general}
Let $n\geq2$ and $\mu\in\Pcal_2(\R)$, and denote the mean of $\mu$ by
$m_\mu$. Then
\begin{equation}\label{eq:general-value}
  \sup_{\pi\in\Pi_n(\mu)}
  \W_2^2(\gamma_n,\pi)
  =
  (n-1)
  +
  \W_2^2\bigl(\gamma,(\sqrt n\,\mathrm{id})_\#\mu\bigr).
\end{equation}
The maximizers are classified as follows.
\begin{enumerate}[label=\textup{(\roman*)}]
  \item If $\mu$ is not symmetric about $m_\mu$, the unique maximizer is
  the diagonal coupling
  \[
    \pi_{\Delta,\mu}
    :=
    \Law(Z,\ldots,Z),
    \qquad Z\sim\mu.
  \]

  \item If $\mu$ is symmetric about $m_\mu$, the maximizers are precisely
  the signed diagonal couplings
  \begin{equation}\label{eq:signed-mu-diagonals}
    \pi_{\varepsilon,\mu}
    :=
    \Law\bigl(
      m_\mu+\varepsilon_1(Z-m_\mu),\ldots,
      m_\mu+\varepsilon_n(Z-m_\mu)
    \bigr),
    \qquad
    \varepsilon\in\{-1,1\}^n,
  \end{equation}
  where $Z\sim\mu$.
\end{enumerate}
\end{theorem}

\begin{proof}[Proof of \Cref{thm:general}]
\emph{Upper bound.}
Fix $\pi\in\Pi_n(\mu)$ and let $T=\nabla\varphi$, where $\varphi \in L^1(\gamma_n)$ is convex, be the Brenier map
from $\gamma_n$ to $\pi$. Note that 
\[
  \mtwo(\gamma_n)=n,
  \qquad
  \mtwo(\pi)=n\,\mtwo(\mu),
\]
and $\nabla\varphi\in L^2(\gamma_n;\R^n)$ because $\| \nabla\varphi\|^2_{L_2(\gamma_n)} =   \mtwo(\pi) < \infty$. Thus Gaussian
integration by parts, \Cref{lem:GaussPartInt}, gives
\begin{align}\label{eq:wasserstein-laplacian}
  \W_2^2(\gamma_n,\pi)
  &=
  n+n\,\mtwo(\mu)
  -2\MC(\gamma_n,\pi) \notag\\
  &=
  n+n\,\mtwo(\mu)
  -2\int_{\R^n} \langle \nabla\varphi(x),x \rangle \, d\gamma_n(x) \notag \\
  &=
  n+n\,\mtwo(\mu)
  -2\int_{\R^n}\rho_n\,d(\Delta\varphi).
\end{align}
As a symmetric positive semidefinite matrix is zero if its trace is zero, $D^2\varphi$ is absolutely continuous w.r.t.\ $\Delta\varphi$. By the Radon--Nikodým theorem, there is a Borel function $A : \R^n \to \R^{n \times n}$ such that
\[
D^2 \varphi = A \Delta \varphi, \qquad \tr(A(x))=1, \qquad A(x) \text{ symmetric positive semidefinite}.
\]
For $T_i=\partial_i\varphi$, we have $T_i\in \BVloc(\R^n)$
and
\[
  D T_i=D^2\varphi\,e_i=(Ae_i)\,\Delta\varphi,
  \qquad
  d|D T_i|=|Ae_i|\,d(\Delta\varphi).
\]
Applying \Cref{lem:matrix} pointwise to $A$ and then
\Cref{prop:ineq_component} to each $T_i$ yields
\begin{align}\label{eq:general-laplacian-bound}
  \int_{\R^n}\rho_n\,d(\Delta\varphi)
  \geq
  \frac1{\sqrt n}\sum_{i=1}^n
  \int_{\R^n}\rho_n\,d|D T_i| 
  \geq
  \frac1{\sqrt n}\sum_{i=1}^n\MC(\gamma,\mu)
  =
  \sqrt n\,\MC(\gamma,\mu).
\end{align}
Combining \eqref{eq:wasserstein-laplacian} and
\eqref{eq:general-laplacian-bound}, we obtain
\[
  \W_2^2(\gamma_n,\pi)
  \leq
  n+n\,\mtwo(\mu)-2\sqrt n\,\MC(\gamma,\mu).
\]
Using that 
\[
  \W_2^2\bigl(\gamma,(\sqrt n\,\mathrm{id})_\#\mu\bigr)
  =
  1+n\,\mtwo(\mu)-2\sqrt n\,\MC(\gamma,\mu),
\]
we conclude 
\[
\W_2^2(\gamma_n,\pi) 
  \leq \W_2^2\bigl(\gamma,(\sqrt n\,\mathrm{id})_\#\mu\bigr) + n-1.
\]

\medskip

\noindent\emph{Construction of maximizers.}
Take $\varepsilon=(1,\ldots,1)$; if $\mu$ is symmetric about
$m_\mu$, let more generally $\varepsilon\in\{-1,1\}^n$. By
\Cref{lem:signed-gradient}, $T_\varepsilon$ is the gradient of a
convex function and is therefore the optimal transport from
$\gamma_n$ to
\[
  \pi_\varepsilon:=(T_\varepsilon)_\#\gamma_n.
\]
Each component of $T_\varepsilon$ has law $\mu$, so
$\pi_\varepsilon\in\Pi_n(\mu)$. If
$\varepsilon=(1,\ldots,1)$, then $\pi_\varepsilon=\pi_{\Delta,\mu}$.
If $\mu$ is symmetric about $m_\mu$, then
\[
  \qmu(\varepsilon_i s)
  =
  m_\mu+\varepsilon_i\bigl(\qmu(s)-m_\mu\bigr)
\]
for almost every $s$, and hence
$\pi_\varepsilon=\pi_{\varepsilon,\mu}$. Moreover,
\begin{align*}
  \int_{\R^n}
  \langle x,T_\varepsilon(x)\rangle\,d\gamma_n(x)
  &=
  \frac1{\sqrt n}\sum_{i=1}^n
  \int_\R \varepsilon_i s\,\qmu(\varepsilon_i s)\,d\gamma(s)\\
  &=
  \sqrt n\int_\R s\,\qmu(s)\,d\gamma(s)
  =
  \sqrt n\,\MC(\gamma,\mu).
\end{align*}
Thus $\MC(\gamma_n,\pi_\varepsilon)
=\sqrt n\,\MC(\gamma,\mu)$, and
\[
  \W_2^2(\gamma_n,\pi_\varepsilon)
  =
  \W_2^2\bigl(\gamma,(\sqrt n\,\mathrm{id})_\#\mu\bigr)+n-1.
\]

\medskip
\noindent\emph{Characterization of equality cases.}
Let $\pi\in\Pi_n(\mu)$ be a maximizer and let $T=\nabla\varphi$
be the optimal map from $\gamma_n$ to $\pi$. If
$\Delta\varphi=0$, then $T$ is constant, so $\mu$ is a Dirac mass
and the conclusion is immediate. We may therefore assume that
$\Delta\varphi\neq0$.

Equality in \eqref{eq:general-laplacian-bound} implies that each
$T_i$ is an equality case in \Cref{prop:ineq_component}. Hence, for
every $i\in\{1,\ldots,n\}$, there is a unit vector $a_i\in\R^n$
such that
\begin{equation}\label{eq:equality-components}
  T_i(x)=\qmu(\langle a_i,x\rangle)
  \qquad\text{for $\gamma_n$-almost every }x.
\end{equation}

Write $B$ for the Radon--Nikodým derivative of $D^2\varphi$ with
respect to $\Delta\varphi$, so that
\[
  D^2\varphi=B\,\Delta\varphi,
  \qquad
  \operatorname{tr}(B)=1,
  \qquad
  B\text{ is symmetric and positive semidefinite}
\]
for $\Delta\varphi$-almost every $x$. It is straightforward to check that $DT_i = a_i ( D\qmu \otimes \mathcal{L}_{a_i^\bot})$, where $\operatorname{span}(a_i)\times a_i^\bot$ is identified with $\R^n$ through $(s a_i,y)\mapsto s a_i+y$ and $\mathcal{L}_{a_i^\bot}$ is Lebesgue measure on $a_i^\bot$. Since $\qmu$ is nondecreasing, $D \qmu \ge 0$ and hence
\[
  DT_i=a_i\,|DT_i|.
\]
On the other hand,
\[
  DT_i=D^2\varphi\,e_i=(Be_i)\,\Delta\varphi,
  \qquad
  |DT_i|=|Be_i|\,\Delta\varphi,
\]
and therefore
\begin{equation}\label{eq:prf:B_struc}
  Be_i=a_i|Be_i|
  \qquad\text{for $\Delta\varphi$-almost every }x.
\end{equation}

Equality in the first inequality of
\eqref{eq:general-laplacian-bound}, together with $\rho_n>0$, yields
\[
  \frac1{\sqrt n}\sum_{i=1}^n|B(x)e_i|=1
\]
for $\Delta\varphi$-almost every $x$. By the equality cases in
\Cref{lem:matrix}, for such $x$ there is
$\sigma(x)\in\{-1,1\}^n$ such that
\[
  B(x)=\frac1n\sigma(x)\sigma(x)^\top.
\]
In particular,
\[
  |B(x)e_i|=\frac1{\sqrt n},
  \qquad
  a_i=\frac{\sigma_i(x)}{\sqrt n}\sigma(x),
\]
where the second identity follows from \eqref{eq:prf:B_struc}.
Taking $i=1$, we see that
\[
  \varepsilon:=\sqrt n\,a_1=\sigma_1(x)\sigma(x)
\]
is a fixed vector in $\{-1,1\}^n$. Consequently, for every $i\in\{1,\ldots,n\}$, 
\[
  a_i=\frac{\varepsilon_i}{\sqrt n}\varepsilon.
\]
It follows from \eqref{eq:equality-components} that
\[
  T_i(x)
  =
  \qmu\left(
    \frac{\varepsilon_i}{\sqrt n}
    \langle\varepsilon,x\rangle
  \right)
  =
  T_{\varepsilon,i}(x).
\]

We may now apply \Cref{lem:signed-gradient}. If all entries of
$\varepsilon$ have the same sign, then $\pi$ is the diagonal
coupling. If $\varepsilon$ contains both signs, the fact that
$T_\varepsilon$ is the gradient of a convex function forces $\mu$ to
be symmetric about $m_\mu$, and then
$\pi=\pi_{\varepsilon,\mu}$. Together with the construction above,
this proves the asserted classification of all maximizers.
\end{proof}

\end{document}
